% Method section (v1, 2026-08-22). Depends on paper/theory.tex (appendix) for the propositions it
% cites. Style: no em-dashes, no prose colons, no numbers, professional register.
% Assumed preamble: amsmath, algorithm/algorithmic (or algorithm2e), natbib.

\section{Method}
\label{sec:method}

We fine-tune a pretrained vision-language-action model on its own demonstrations, entirely offline,
so that it improves along two dimensions at once: the actions inside a chunk, and how much of that
chunk to execute before querying again. The method has three parts. A critic assigns a value to
every prefix of a chunk, trained so that commitments of different lengths are comparable
(Section~\ref{sec:method:critic}). The policy is improved against that critic on the full chunk
(Section~\ref{sec:method:actor}). At deployment the same critic chooses how long to commit, with no
sampling and no ranking of candidates (Section~\ref{sec:method:deploy}).

\subsection{Setup and notation}
\label{sec:method:setup}

A pretrained policy $\pi_\theta$ maps an observation $o$ to a chunk of $H$ actions. Writing
$c=a_{1:H}$ for a chunk and $c_{1:k}$ for its first $k$ actions, executing $k$ actions from state
$s$ and then querying the policy again is a temporally extended action of duration $k$. We keep the
backbone frozen and attach a critic to its representations, so fine-tuning changes the action
head and the commitment rule rather than the perceptual model.

Because commitments have different durations, values must be compared with semi-Markov
bookkeeping~\citep{sutton1999options, bradtke1994smdp},
\begin{equation}
\label{eq:method-smdp}
Q^{\pi,\kappa}(x, c, k) \;=\; \mathbb E\Big[\textstyle\sum_{j<k}\gamma^{j} r_j \;+\; \gamma^{k}\,
V^{\pi,\kappa}(x_k)\Big],
\end{equation}
where $\kappa$ is the commitment rule and $x$ denotes the critic's input, defined next. Charging a
single discount per decision instead of per elapsed step distorts values by a term that grows with
$k$ and whose sign follows the reward convention, so the induced preference over commitments is an
artifact of where the value scale places its zero rather than a property of the task
(Proposition~\ref{prop:bookkeeping} and Remark~\ref{rem:sign}).

\subsection{A critic that prices commitment fairly}
\label{sec:method:critic}

Three properties of demonstration data make a naive chunk critic misprice commitment, and each is
addressed by a specific choice rather than by a loss term.

\paragraph{History as input.}
Human demonstrations are not Markovian in the observation. A teleoperator holds a plan across a
stretch of motion and acts from it while the scene alone no longer reveals it, which is exactly what
makes a committed execution informative. A critic conditioned on the instantaneous observation has
no room to represent that information, so it cannot distinguish continuing a plan from resuming
one. We therefore condition on $x_t=(o_t, a_{t-m:t-1})$, the observation together with the actions
already executed since the last query. This is a change of input, not of objective, and it is what
Remark~\ref{rem:confound} requires when the latent lies in the past.

\paragraph{Interventional composition.}
The value of a $k$-step commitment must not be estimated by regressing outcomes on the executed
chunk. In closed-loop demonstrations, an event revealed inside the window causes both the recorded
actions and the outcome, so such a regression conditions on the demonstrator's foreknowledge and
credits the open-loop executor with information it will not have. Composing one-step backups is
necessary but not sufficient, since bootstrapping each transition from its own recorded successor
preserves the same correlation. We therefore compose the backup while resampling the successor
among transitions at the same decision point, holding the evaluated tail fixed,
\begin{equation}
\label{eq:method-interv}
y_k(x_t, c) \;=\; r_t \;+\; \gamma\, \mathbb E_{x' \sim \mathcal{D}(\cdot \mid x_t)}
\big[\, Q_{\bar\phi}(x', c_{2:k}, k-1) \,\big],
\qquad k > 1,
\end{equation}
where $\mathcal D(\cdot\mid x_t)$ is the empirical distribution of successors at decision points
matching $x_t$ and $\bar\phi$ are target parameters. The target then estimates
$Q(x,\mathrm{do}(c_{1:k}))$, which is what open-loop execution faces. Under open-loop consistent
data the two coincide~\citep{dqc}; the intervention is what buys that independence when the data is
closed loop.

\paragraph{An honestly priced requery.}
The leaf of the recursion is the value of querying the policy again, and it must be the value of
the policy that will actually resume control. Bootstrapping from demonstration returns instead
credits the requery with the demonstrator's competence, an optimism that enters with weight
$\gamma^{k}$ and therefore biases the choice toward short commitments regardless of the reward
convention (Proposition~\ref{prop:optimism}). We use the deployed policy's own expectation,
\begin{equation}
\label{eq:method-leaf}
y_1(x_t, c) \;=\; r_t \;+\; \gamma\, \mathbb E_{x'}\,\mathbb E_{c'\sim\pi_\theta(\cdot\mid o')}
\big[\, Q_{\bar\phi}\big(x', c', \kappa(x', c')\big) \,\big],
\end{equation}
so the critic is a fixed point of the process that will run, not of an idealized one. Targets stay
within the support of the data in the sense that the chunks scored during training are the recorded
ones; only the leaf involves the current policy, one step deeper, where the state is more
informative.

The critic is a causal transformer over the chunk on frozen backbone features, emitting all $H$
prefix values in one pass, and is trained by regression on
$\{y_k\}_{k=1}^{H}$ from \eqref{eq:method-interv} and \eqref{eq:method-leaf} with a target network.
The state baseline used below shares the trunk with no head of its own only if an ablation supports
it; we treat this as a measured choice rather than a design commitment.

\subsection{Improvement on the full chunk}
\label{sec:method:actor}

Selection cannot lengthen commitments. If the policy is frozen and only the commitment rule is
tuned, the option values are fixed and the deployed behavior is a deterministic functional of the
initial policy; and improving only the selected short prefixes leaves the full-chunk value
unchanged, which is the quantity that lower bounds deployed value
(Lemma~\ref{lem:fullchunk} and Remark~\ref{rem:frozen}). We therefore improve the policy against
the same critic on the \emph{full} chunk, by advantage-weighted regression on the demonstrations,
\begin{equation}
\label{eq:method-actor}
\max_\theta\; \mathbb E_{(x, c)\sim\mathcal D}\Big[\,
w(x,c)\,\log \pi_\theta(c \mid o)\Big],
\qquad
w(x,c)=\exp\!\big(\beta^{-1}\big[Q_\phi(x, c, H)-V_\phi(x)\big]\big),
\end{equation}
with $V_\phi(x)=\mathbb E_{c'\sim\pi_\theta}[Q_\phi(x,c',H)]$. Weighting the recorded chunks keeps
the update inside the demonstrations, so the critic is never maximized over chunks it has not
scored, which is the failure mode that makes deterministic policy-gradient extraction unstable
offline. Nothing here is sampled and ranked at deployment; the sampling that occurs during training
is compiled into the policy.

\subsection{Deployment}
\label{sec:method:deploy}

At each query the policy emits one chunk, the critic scores its $H$ prefixes in the same forward
pass, and the system executes the longest prefix whose value is within $\epsilon$ of the best,
\begin{equation}
\label{eq:method-kappa}
\kappa(x, c) \;=\; \max\Big\{\,k \;:\; Q_\phi(x, c, k) \;\ge\; \max_{k'} Q_\phi(x, c, k') - \epsilon
\,\Big\}.
\end{equation}
The tie-breaking direction matters. Values of neighbouring commitments are often
indistinguishable, and a rule that breaks ties toward short commitments requeries constantly, which
re-injects policy error and destroys the plan continuity that made the demonstration succeed;
breaking them toward long commitments yields the curriculum that improvement predicts, as the set
of states where a strictly shorter commitment wins shrinks to those where reactive information
genuinely arrives (Corollary~\ref{cor:curriculum}). The comparison tolerance $\epsilon$ is the only
free parameter introduced by the commitment rule, and no cost term is added to the reward.

\subsection{Algorithm}
\label{sec:method:algo}

Training alternates critic regression and policy improvement on a fixed offline dataset. One pass
is: fit the critic to \eqref{eq:method-interv} and \eqref{eq:method-leaf} against the current
policy; improve the policy with \eqref{eq:method-actor}; repeat. Only the action head and the
critic receive gradients; the backbone stays frozen throughout. No environment interaction is used
at any point, and the deployed system performs exactly one policy query and one critic pass per
decision.
